%%%%%%%%%%%%%%%%%%%%%%%%%%%%%%%%%%%%%%%%%%%%%%%%%%%
%%% ESTADO DEL ARTE
%%%%%%%%%%%%%%%%%%%%%%%%%%%%%%%%%%%%%%%%%%%%%%%%%%%

\chapter{Estado del arte}
\fancyhead[RE]{\textsc{CAPÍTULO} \thechapter. Estado del arte}
\label{ch:EstadoArte}


\section{Historia del problema}
\label{sec:JobTitleMatching}

El \gls{JTM} constituye una tarea central dentro del campo más amplio de la gestión del talento y la \gls{HCM}. En su forma básica, el problema consiste en determinar la equivalencia o el grado de compatibilidad entre títulos de puestos, a pesar de la enorme variabilidad léxica, la polisemia y las convenciones internas de las organizaciones. Esta variabilidad hace que una misma ocupación pueda aparecer bajo campos muy distintos y, a su vez, que títulos idénticos traduzcan responsabilidades distintas según el contexto organizativo, sectorial o geográfico.

\subsection{Métodos tradicionales y enfoques económicos}

Históricamente, los trabajos sobre emparejamiento laboral han adoptado perspectivas tanto técnicas como económicas. En economía laboral española existe una persistencia en el proceso de emparejamiento que muestra cómo las fricciones de mercado y la heterogeneidad ocupacional condicionan los resultados del matching laboral \cite{de2008funcion}. Esos estudios clásicos proporcionan un marco para interpretar por qué la normalización de titulaciones no es sólo un problema lingüístico, sino también institucional y estructural: la misma denominación puede implicar distintas funciones y niveles de especialización en diferentes segmentos del mercado laboral \cite{de2017quien}.

Desde el punto de vista computacional, los primeros abordajes aplicados al dominio laboral siguieron estructuras generales de recuperación de información y clasificación de texto: normalización de cadenas, reglas heurísticas y distintas métricas de similitud. Aunque estas técnicas (y los sistemas basados en reglas) pueden ser útiles en entornos controlados, su coste de mantenimiento y su fragilidad frente a sinónimos, abreviaturas y errores hacen que resulten insuficientes para datos del mundo real y escalas grandes \cite{jirjees2025machine}.

\subsection{Transición hacia técnicas de aprendizaje automático}

La llegada y consolidación del \gls{ML} transformó la investigación y la práctica en \gls{JTM}. Durante los últimos años, la literatura documenta un amplio abanico de técnicas: desde modelos supervisados clásicos hasta arquitecturas neuronales que aprenden representaciones textuales más ricas y robustas frente al ruido y la variabilidad \cite{jirjees2025machine}. Han comenzado a surgir soluciones no supervisadas y semi-supervisadas para detectar patrones, agrupar titulaciones y reducir la dependencia de grandes volúmenes de anotación manual \cite{mamani2022machine,mamani2022proceso}.

Específicamente en entornos hispanohablantes, trabajos recientes han explorado procesos de clustering y pipelines de \gls{ML} no supervisado para identificar la demanda social de puestos y determinar similitudes entre títulos de profesionales de sectores específicos. Estas contribuciones muestran que, aunque las técnicas de clustering capturan agrupaciones útiles, su sensibilidad a la representación textual y su capacidad de generalización condiciona fuertemente la calidad del emparejamiento.

En cuanto a los análisis bibliométricos sobre investigación en esta tarea revelan un crecimiento sostenido del interés científico en la intersección entre \gls{AI} y reclutamiento por parte de la \gls{HRM}, con aplicaciones que abarcan desde selección automatizada hasta orientación profesional y análisis de brechas de competencias \cite{rojas2022bibliometric}. Dichos trabajos también evidencian que, aunque existen esfuerzos notables en sectores como la salud o la tecnología, persiste una desproporción entre la disponibilidad de recursos en inglés y en otras lenguas, lo que limita la comparabilidad y la replicabilidad de soluciones en contextos locales.

\subsection{Del problema de clasificación a la tarea de \emph{matching} en entornos reales}

Una distinción clave que ha emergido en la literatura reciente es la diferencia entre (i) la clasificación contra una taxonomía canónica y (ii) el matching o ranking entre pares de textos libres (oferta vs. perfil solicitante). Mientras la clasificación se beneficia de puntos de referencia estables, el matching real debe lidiar con ruido, abreviaturas, títulos internos y asimetrías informativas que no aparecen en taxonomías limpias. En consecuencia, las soluciones más robustas tienden a combinar representaciones semánticas aprendidas con estrategias de evaluación y depuración específicas del dominio, integrando tanto medidas globales de similitud como señales estructurales extraídas del mercado laboral \cite{gasco2025brief}.

%%

\section{Enfoques abordados en TalentCLEF}
\label{sec:TalentCLEFOverview}

La Tarea A de TalentCLEF 2025, centrada en el Multilingual \gls{JTM}, atrajo una variedad de enfoques que buscaron resolver el desafío de alinear semánticamente títulos de trabajo en inglés, español, alemán y, opcionalmente, chino. El análisis de las soluciones presentadas por los participantes revela una fuerte tendencia hacia el uso de modelos de lenguaje basados en Transformers, aunque las estrategias de implementación, entrenamiento y uso de recursos externos variaron significativamente.

\subsection{Estrategias Dominantes}

La mayoría de las soluciones giraron en torno a modelos de \textit{Sentence Transformers} multilingües preentrenados, como MPNet \cite{brikman2025multilingual}, JobBERT \cite{decorte2025multilingual} o variantes de SBERT \cite{nizami2025enhancing,bhangale2025fine}. El objetivo común fue la generación de embeddings en un espacio semántico compartido, donde títulos de trabajo equivalentes en diferentes idiomas tuvieran una alta similitud\footnote{La similitud medida en los espacios vectoriales generados fue la similitud coseno, elegida de forma unánime entre todas las propuestas de solución al problema.}.

Un hallazgo clave, destacado por el equipo NLPnorth \cite{zhang2025nlpnorth}, fue la comparación directa de tres métodos de ajuste fino:
\begin{itemize}
    \item \textbf{Discriminativo (Clasificación):} Reformular la tarea como una clasificación binaria (si dos títulos coinciden o no), generando para ello pares negativos.
    \item \textbf{Contrastivo:} Entrenar el modelo para acercar las representaciones de pares positivos y alejar las de pares negativos.
    \item \textbf{Basado en Prompts:} Utilizar LLMs para evaluar la similitud.
\end{itemize}
Los resultados de NLPnorth sugirieron que el aprendizaje contrastivo fue su estrategia más efectiva.

Otros equipos, como UDII-UPM \cite{rodriguez2025udii}, exploraron el rendimiento de modelos preentrenados sin ajuste fino, aplicando un enfoque \textit{zero-shot}, proponiendo una fusión de representaciones semánticas de diferentes modelos para robustecer el matching. Por el contrario, el equipo Pjmathematician utilizó embeddings GIST e incorporó datos generados mediante LLMs \cite{vachharajani2025pjmathematician}.

\subsection{El Rol Crítico de Recursos Externos}

El uso de la taxonomía \gls{ESCO} fue un diferenciador clave para varias de las soluciones de mayor rendimiento.

\subsubsection{Uso de ESCO}
La taxonomía \gls{ESCO} no solo sirvió como base de conocimiento, sino también como fuente de datos de entrenamiento. El equipo VerbaNex \cite{novoa2025verbanex} utilizó las relaciones jerárquicas de ESCO para generar automáticamente pares adicionales positivos y negativos, ajustando su modelo para crear un espacio semántico alineado con la taxonomía. De manera similar, NLPnorth \cite{zhang2025nlpnorth} también extrajo datos adicionales de ESCO para su entrenamiento\footnote{Cabe destacar que el conjunto de datos de entrenamiento ya habían sido generado por la organización y se encontraba a disposición de los participantes. Sin embargo, conociendo la procedencia y método de producción de este set de datos, algunos participantes decidieron aumentar dicho set para una mejor generalización de los modelos en el aprendizaje.}.

Otros equipos adoptaron un enfoque híbrido. El sistema de dos etapas de Rodríguez et al. \cite{rodriguez2025two} utilizó \gls{ESCO} como un primer filtro o sistema de recuperación de candidatos, para luego refinar la lista de candidatos mediante un LLM en una etapa de \textit{re-ranking}.

\subsubsection{Uso de LLMs}
Además del \textit{re-ranking} mencionado, varios equipos utilizaron algún \gls{LLM} para la generación de datos. El equipo NT \cite{nga2025nt} propuso un sistema basado en la generación de descripciones de trabajo mediante LLMs, seguida de una recuperación de información. Sin embargo, su rendimiento fue bastante limitado en comparación con otros enfoques basados en Transformers.

\subsection{Gestión de Escenarios Multilingües y Cross-Lingües}

El desafío principal de la tarea era la correcta alineación semántica entre idiomas.
\begin{itemize}
    \item \textbf{Espacio Semántico Compartido:} La estrategia dominante fue el uso de codificadores multilingües (SBERT, MPNet) ajustados para proyectar todos los idiomas en un único espacio vectorial \cite{brikman2025multilingual,nizami2025enhancing,novoa2025verbanex}.
    \item \textbf{Inglés como Lengua Pivote:} El equipo de Rodríguez et al. \cite{rodriguez2025two} investigó explícitamente el uso del inglés como idioma pivote, traduciendo los títulos de otros idiomas al inglés antes de la comparación, una estrategia que demostró mejorar el rendimiento en algunos escenarios cross-lingües.
    \item \textbf{Recuperación Híbrida:} El equipo AlexU-NLP \cite{barakat2025alexu} propuso un sistema de recuperación híbrido con un enfoque de \textit{curriculum learning}, diseñado para gestionar la dificultad progresiva de los pares multilingües.
\end{itemize}

\subsection{Consideraciones Específicas: Chino y Sesgo de Género}

Aunque la mayoría de los participantes se centraron en inglés, español y alemán, la tarea incluía dos consideraciones importantes: el idioma chino y la evaluación del sesgo de género.

\subsubsection{El Caso del Chino}
El idioma chino (en pares \textit{zh-zh} y \textit{en-zh}) fue una pista opcional dentro de la Tarea A. Sin embargo, la mayoría de los trabajos centraron sus esfuerzos en los idiomas europeos, debido a que el chino no se encontraba dentro de los datos de entrenamiento, sino únicamente en validación y testeo, con la intención de evaluar la capacidad multilingue de las distintas propuestas. 

\subsubsection{Evaluación del Sesgo de Género}
Una componente evaluativa clave de la tarea fue la medición del sesgo de género. En idiomas como el español y el alemán, donde los títulos de trabajo pueden tener marcas explícitas de género, se evaluó cómo las distintas soluciones manejaban estas diferencias,  para evaluar si los sistemas perpetuaban o mitigaban sesgos de género en sus predicciones. Los datos de la competición incluían anotaciones ocultas de género y expresiones con marcas de género para permitir este análisis. La organización proponía buscar un sistema que minimizara las diferencias de rendimiento entre títulos de trabajo asociados a diferentes grupos de género.

A pesar de que esta evaluación fue central, pocos trabajos de los participantes detallaron explícitamente las técnicas de mitigación del sesgo que aplicaron. La mayoría de los sistemas se centraron en la precisión del matching semántico. Los trabajos de NLPnorth y de Ixa\todo{Revisar nombre de equipos y de autores, para que haya unificación en la forma en la que los llamo.} evaluarion este tema con detenimiento, consiguiendo una puntuación perfecta en la métrica de sesgo de género.

%%

\section{Transformers para Semejanza Semántica}
\label{sec:Transformers}

Los transformers\todo{No se si todas las palabras del inglés deben ponerse en cursiva o no hace falta.} han sido la palanca tecnológica dominante en las tareas de comprensión y comparación semántica de textos desde su aparición. En el contexto de la tarea, su adopción no sólo ha permitido avances significativos en calidad de ranking, sino que ha generado un conjunto diverso de arquitecturas y estrategias de uso que conviene desglosar.

\subsection{Definición y características}
Un transformer es una arquitectura de red neuronal basada principalmente en mecanismos de atención que procesan secuencias sin recurrir a operaciones recurrentes o convolucionales. La atención permite que cada posición de entrada sintetice información de todas las demás posiciones mediante pesos aprendidos; esto facilita la captura de relaciones de largo alcance, la paralelización durante el entrenamiento y la construcción de representaciones contextualizadas de tokens y oraciones \cite{vaswani2017attention}. Aunque el diseño original incluye bloques de encoder y decoder, la práctica moderna para tareas de recuperación y matching usa variantes sólo-encoder (para embeddings), encoder-decoder y modelos sólo-decoder (para \gls{LLM}) según el uso concreto.

\subsection{Su utilización en la tarea}
En la práctica aparecen tres patrones arquitectónicos generales, con variantes e intermedios frecuentes:

\subsubsection{Bi-encoders (sentence-transformers / dual encoders)}
Cada query se codifica por separado con el mismo encoder para producir vectores en un espacio semántico compartido, y usa la similitud entre vectores para ordenar candidatos. Este esquema es ideal cuando se requiere recuperar a escala porque la indexación de vectores y la búsqueda aproximada permiten latencias bajas. Su principal fortaleza es la eficiencia en la fase de inferencia; su principal limitación es que la interacción entre consulta y candidato es sólo implícita en el espacio vectorial, por lo que puede fallar a la hora de resolver matices finos que requerirían un cotejo token-a-token.

\subsubsection{Cross-encoders (re-rankers y matchers de interacción)}
Un cross-encoder toma la pareja (consulta, candidato) y las procesa juntas, produciendo una puntuación que refleja la compatibilidad tras un proceso de atención en ambas direcciones. Proporciona decisiones más precisas porque permite interacciones finas entre las dos secuencias, capturando coincidencias léxicas localizadas, supresiones y dependencias sintácticas que un bi-encoder puede pasar por alto. Sin embargo, el coste computacional es mucho más alto. Con esta técnica no es viable aplicar cross-encoders exhaustivamente sobre grandes conjuntos sin una fase previa de filtrado. Por ello se emplean habitualmente como re-rankers sobre un conjunto reducido (top-k) obtenido con otros métodos menos costosos. La práctica de combinar un retrieval rápido y re-ranking posterior ha quedado consolidada como patrón de diseño en sistemas de \gls{IR} modernos \cite{sbertRetrieveamp,xu2025survey}.

\subsubsection{Modelos autoregresivos y LLMs para generación y scoring}
Los modelos autoregresivos se han usado de forma complementaria: para generar descripciones adicionales de un título, traducir o normalizar textos, crear negativos/positivos sintéticos, o para puntuar directamente parejas mediante prompting. Ofrecen flexibilidad y poder de razonamiento, pero también son muy propensos a la generación de ruido. Requieren de un control cuidadoso y suelen tener costes computacionales y de infraestructura elevados.

\subsection{Técnicas de entrenamiento específicas para mejorar la calidad del ranking}
El rendimiento de los transformers depende tanto de la arquitectura como de la estrategia de entrenamiento y de la calidad de los datos usados. En la tarea se observan varias técnicas que marcan la diferencia:

\begin{itemize}
    \item \textbf{Entrenamiento contrastivo:} las técnicas contrastivas, como por ejemplo \gls{InfoNCE}, acercan los pares positivos de vectores y separan los negativos. Una selección inteligente de negativos mejora la discriminación del espacio semántico y reduce falsos positivos en los rankings. Técnicas recientes que centran su foco en la selección de negativos han mostrado ganancias notables en el fine-tuning de embeddings \cite{solatorio2024gistembed}.
    \item \textbf{Variantes de las técnicas contrastivas:} existen algunas modificaciones, como focal-\gls{InfoNCE}, que tratan de penalizar más los errores en ejemplos difíciles y estabilizan el aprendizaje cuando existe un desbalance fuerte entre positivos y negativos. Con ello consiguen mejorar la capacidad de priorizar los verdaderos positivos en los primeros puestos del ranking \cite{hou2023improving}.
    \item \textbf{Curriculum learning:} comenzar con negativos fáciles y evolucionar hacia negativos más difíciles (o introducir negativos guiados por similitud léxica/semántica) puede estabilizar el entrenamiento y evitar que el modelo colapse ante una distribución ``long-tail'' de títulos.
    \item \textbf{Fine-tuning multilingüe:} ajustar modelos multilingües con datos balanceados por idioma puede ayudar a alinear representaciones cross-lingües y a reducir la fragmentación de la representación del embedding. Algunos métodos modernos de entrenamiento masivo multilingüe han demostrado que el aprendizaje dirigido para varios idiomas simultáneamente mejora la transferencia entre pares idiomáticos \cite{mao2024ems}.
\end{itemize}\todo{No se si debo entrar más en detalle sobre InfoNCE o técnicas concretas que vaya a usar en el modelo, o si las debo explicar luego en otra sección.}

% \subsection{Beneficios de unas técnicas sobre otras}
% La elección entre bi-encoders, cross-encoders y LLM-augmented pipelines obedece a un trade-off clásico entre \emph{eficiencia} y \emph{precisión}:

% \begin{itemize}
%     \item \textbf{Bi-encoders:} su mayor ventaja es la escalabilidad y la latencia baja en sistemas de producción; son excelentes para obtener candidatos iniciales. Funcionan muy bien cuando se ha invertido en un fine-tuning cuidadoso con negatives difíciles.
%     \item \textbf{Cross-encoders:} imprescindibles cuando la tarea exige distinciones semánticas finas o cuando el coste de equivocarse en el top-1 es elevado. Su uso como re-ranker combina lo mejor de ambos mundos: velocidad en el retrieval y precisión en la selección final.
%     \item \textbf{LLMs y generación:} permiten enriquecer los ejemplos y crear señales que los encoders tradicionales no capturan (explicaciones, desambiguaciones, traducciones). Útiles para datos escasos o para lenguas con poca anotación directa. Sin embargo, introducir generación automática requiere controles de calidad y mecanismos para evitar que el ruido sintético degrade los embeddings.
% \end{itemize}

% \subsection{Desventajas y riesgos prácticos}
% A pesar de su eficacia, los transformers tienen limitaciones que condicionan su empleo en aplicaciones de \gls{HCM}:

% \begin{itemize}
%     \item \textbf{Coste computacional y latencia:} los cross-encoders y los LLMs son caros en tiempo y memoria. En producción es necesario combinar técnicas de distillation, early-exit o re-ranking para mantener aceptables latencias.
%     \item \textbf{Necesidad de datos y sensibilidad al dominio:} los transformers preentrenados requieren fine-tuning con datos representativos del dominio (títulos laborales, abreviaturas sectoriales). El dominio laboral es heterogéneo; sin ejemplos suficientes, el rendimiento puede degradarse por shift de dominio.
%     \item \textbf{Tokenización y particularidades lingüísticas:} las decisiones de tokenización (subword, caracteres, vocabularios multilingües) afectan el tratamiento de idiomas como el chino, donde la segmentación palabra/ carácter influye en la calidad de las representaciones; estudios recientes han mostrado que combinar granularidades mejora el comportamiento de modelos preentrenados en chino \cite{liang2023character}.
%     \item \textbf{Sesgos y equidad:} los transformers reflejan sesgos presentes en los datos de preentrenamiento y fine-tuning. En el contexto del emparejamiento laboral, esto puede traducirse en diferencias de ranking por género, edad o nacionalidad si no se evalúa y mitiga explícitamente; la literatura reciente ofrece marcos y técnicas para detección y mitigación que conviene incorporar al ciclo de evaluación \cite{nemani2024gender}.
%     \item \textbf{Robustez frente al ruido y abreviaturas:} aunque las representaciones contextualizadas ayudan con sinonimia, los títulos internos y formatos erróneos pueden seguir confundiendo los modelos sin un preprocesamiento o estrategias de augmentación específicas.
% \end{itemize}

% \subsection{Consideraciones de implementación}
% Para integrar transformers en un sistema de ranking multilingüe eficaz conviene atender a varios aspectos prácticos:

% \begin{itemize}
%     \item \textbf{Infraestructura de retrieval:} indexación de vectores y búsqueda ANN (approximate nearest neighbor) para bi-encoders; pipeline de recuperación + re-ranking para equilibrar costo y precisión \cite{sbertRetrieveamp}.
%     \item \textbf{Estrategia de evaluación:} además de métricas tradicionales (MAP, MRR), usar medidas de estabilidad de ranking y pruebas por subgrupos (género, idioma) para detectar sesgos y verificación de calidad por idioma.
%     \item \textbf{Ciclo de vida del modelo:} monitorizar deriva de dominio, actualizar con nuevos ejemplos reales y mantener un proceso de curado/filtrado de ejemplos sintéticos provenientes de LLMs.
%     \item \textbf{Eficiencia y compactación:} técnicas de distillation, quantization y early-exit permiten desplegar modelos de menor coste manteniendo buena parte del rendimiento.
% \end{itemize}

% \subsection{Resumen}
% Los transformers ofrecen un marco potente y flexible para Task A: permiten aprender representaciones semánticas robustas, facilitar la transferencia cross-lingüe y sostener pipelines híbridos que combinan rapidez y precisión. No obstante, su efectividad depende de prácticas cuidadosas de entrenamiento (negative mining, curriculum, augmentación), decisiones de tokenización en idiomas específicos como el chino, y de una evaluación amplia que incorpore aspectos de equidad y robustez. La arquitectura óptima en TalentCLEF —como se ha visto entre los participantes— suele ser una combinación pragmática: bi-encoder afinado para retrieval, generación/augmentación controlada con LLMs para enriquecer datos, y cross-encoder o LLM-based re-ranking para maximizar la precisión en el top-k.
